% !TEX root = ../../main.tex

\section{Alcance y limitaciones de la validación}
\label{sec:validacion-alcance}

Una vez explicada la metodología de validación y presentados los programas sobre los que se aplica, es necesario delimitar el alcance de la evidencia obtenida y explicitar las limitaciones del proceso. Estas consideraciones permiten interpretar adecuadamente los resultados presentados en el Capítulo~\ref{chap:resultados} y evitar conclusiones que excedan los casos evaluados.

Primeramente, el alcance de la evaluación está condicionado por la naturaleza sintética de los 99 programas que componen los corpus. Estos fueron construidos para aislar de manera controlada los fenómenos relevantes para la solución, pero no permiten estimar con qué frecuencia aparecen oportunidades de reordenamiento en aplicaciones reales. La cobertura tampoco abarca toda la sintaxis de \textsf{Haskell}: en particular, la implementación no incorpora lecturas que aparezcan exclusivamente dentro del patrón de un \texttt{BindStmt}, como puede ocurrir con \texttt{ViewPatterns}. De manera similar, los controles de activación consideran las tres configuraciones definidas en la metodología, pero no todas las combinaciones posibles de extensiones, importaciones y formas de calificar un bloque \texttt{do}.

Además de esta cobertura acotada, la evidencia asociada a \textbf{PV1} es de carácter observacional. La comparación de la salida capturada y del código de salida permite detectar diferencias en las ejecuciones de programas cerrados y finitos, pero no constituye una demostración formal de equivalencia semántica ni cubre divergencia o interacciones externas. En los programas probabilísticos, \texttt{Dist.norm} agrupa resultados iguales y suma sus probabilidades antes de imprimir. Aunque esta normalización evita confundir diferencias de representación interna con diferencias semánticas, también pasa a formar parte del criterio observable adoptado.

Incluso cuando las ejecuciones producen resultados coincidentes, su interpretación depende de la conmutatividad declarada por el programador, pues \textsf{GHC} no demuestra esta propiedad. Si la declaración es incorrecta, preservar las dependencias RAW no basta para conservar el comportamiento. El control con \texttt{IO} muestra este límite: un reordenamiento puede ser estructuralmente válido y, aun así, modificar el orden de los efectos observables. Por tanto, los resultados positivos de \texttt{Maybe} y \texttt{Dist.T Rational} deben entenderse bajo la hipótesis explícita de conmutatividad adoptada por la memoria.

En cuanto a \textbf{PV2}, existe una limitación adicional porque las métricas, los candidatos y las decisiones provienen del mismo prototipo que se evalúa. Para evitar que el resumen interno provisto por el Renamer sea el único criterio de aceptación, el flujo compila cada candidato por separado, vuelve a calcular el mínimo y contrasta cantidades de reordenamientos que pueden determinarse independientemente. Aun así, no existe un oráculo externo que verifique todos los grafos y costos producidos.

Finalmente, aun cuando la selección sea correcta en los casos evaluados, la extensión enfrenta límites de escalabilidad y de interpretación del costo. La enumeración exhaustiva de ordenamientos topológicos puede crecer factorialmente en el peor caso, especialmente cuando todos los statements son independientes, por lo que su aplicación puede resultar impracticable en bloques sumamente grandes. A su vez, el modelo asigna un costo uniforme a cada statement y compara únicamente la estructura de los planes; no distingue operaciones baratas de costosas, no incorpora la sobrecarga de los combinadores ni mide paralelismo efectivo. En consecuencia, la evidencia permite validar la selección de un mínimo estructural bajo el modelo adoptado, pero no garantiza la escalabilidad de la extensión ni una mejora de rendimiento en tiempo de ejecución.
